% The opening pages of G. H. Hardy, Divergent Series (Oxford, Clarendon
% Press, 1949), chapter I, pp. 1-6, for comparing against the printed book.

\chapter{Introduction}

\section{The sum of a series.} The series
$$\sum_0^\infty a_n = a_0+a_1+a_2+\dots$$
is said to be \emph{convergent}, to the sum $s$, if the `partial sum'
$$s_n = a_0+a_1+\dots+a_n$$
tends to a finite limit $s$ when $n\to\infty$; and a series which is not
convergent is said to be \emph{divergent}. Thus the series
$$\no\label{1-1} 1-1+1-1+\dots, \no 1-2+3-4+\dots, \\
  \no 1-2+4-8+\dots, \no 1-1!+2!-3!+\dots, \\
  \no 1+1+1+1+\dots, \no 1+2+4+8+\dots,$$
are divergent. The series
$$\no 1+e^{i\theta}+e^{2i\theta}+\dots, \no \frac12+\cos\theta+\cos2\theta+\dots,$$
are divergent for all real $\theta$, and
$$\no \sin\theta+\sin2\theta+\sin3\theta+\dots$$
is divergent except when $\theta$ is a multiple of $\pi$, when it converges to
the sum 0.

The definitions of convergence and divergence are now commonplaces of
elementary analysis. The ideas were familiar to mathematicians before Newton
and Leibniz (indeed to Archimedes); and all the great mathematicians of the
seventeenth and eighteenth centuries, however recklessly they may seem to
have manipulated series, knew well enough whether the series which they used
were convergent. But it was not until the time of Cauchy that the
definitions were formulated generally and explicitly.

Newton and Leibniz, the first mathematicians to use infinite series
systematically, had little temptation to use divergent series (though
Leibniz played with them occasionally). The temptation became greater as
analysis widened, and it was soon found that they were useful, and that
operations performed on them uncritically often led to important results
which could be verified independently. We give a few simple examples in the
next section; in Ch.~II we shall give others, of greater importance, from
the work of the classical analysts.

\section{Some calculations with divergent series.} We know that
$$\no\label{geom} 1+x+x^2+\dots = \frac{1}{1-x}$$
if $|x|<1$. It seems plain that, if we are to attribute a `sum', in some
sense, to the series for other $x$, this sum should be formally the same.
For (i) it would be very inconvenient if the formula varied in different
cases; (ii) we should expect the sum $s$ to satisfy the equations
$$s = 1+x+x^2+x^3+\dots = 1+x(1+x+x^2+\dots) = 1+xs;$$
and (iii) the left-hand side of \eqref{geom} is the result of performing the
division implied by the right, so that there is certainly one sense of
`$=$' with which \eqref{geom} may be said to be true for all $x$.

(1) Let us assume then that \eqref{geom} is, in some sense, true for all $x$
(except perhaps for $x=1$, which plainly presents special difficulties), and
operate on the formula in an entirely uncritical spirit.

Putting $x=e^{i\theta}$, where $0<\theta<2\pi$ (so that $x\ne1$), we obtain
$$\no 1+e^{i\theta}+e^{2i\theta}+\dots = (1-e^{i\theta})^{-1} = \frac12+\frac12i\cot\frac12\theta,$$
and so
$$\no \frac12+\cos\theta+\cos2\theta+\dots = 0, \no\label{sin} \sin\theta+\sin2\theta+\dots = \frac12\cot\frac12\theta,$$
for $0<\theta<2\pi$. Changing $\theta$ into $\theta+\pi$, we obtain
$$\no\label{cos-alt} \frac12-\cos\theta+\cos2\theta-\dots = 0, \no\label{sin-alt} \sin\theta-\sin2\theta+\dots = \frac12\tan\frac12\theta,$$
for $-\pi<\theta<\pi$. For $\theta=0$, \eqref{cos-alt} gives
$$\no\label{grandi} 1-1+1-\dots = \frac12.$$

(2) We now differentiate \eqref{cos-alt} and \eqref{sin-alt} repeatedly with
respect to $\theta$. We thus obtain
$$\no\label{d1} \sum_1^\infty (-1)^{n-1}n^{2k}\cos n\theta = 0 \quad (k=1,2,\dots;\ -\pi<\theta<\pi), \\
  \no\label{d2} \sum_1^\infty (-1)^{n-1}n^{2k+1}\sin n\theta = 0, \\
  \no \sum_1^\infty (-1)^{n-1}n^{2k}\sin n\theta = (-1)^k\left(\frac{d}{d\theta}\right)^{2k}\frac12\tan\frac12\theta, \\
  \no\label{d4} \sum_1^\infty (-1)^{n-1}n^{2k+1}\cos n\theta = (-1)^k\left(\frac{d}{d\theta}\right)^{2k+1}\frac12\tan\frac12\theta,$$
the last three formulae for $k=0$, $1,\dots$, $-\pi<\theta<\pi$. In
particular, putting $\theta=0$ in \eqref{d1} and \eqref{d4}, and
$\theta=\frac12\pi$ in \eqref{d2}, and remembering that the Taylor's series
for $\frac12\tan\frac12\theta$ is
$$\frac12\tan\frac12\theta = \Sum_{k=0}^\infty \frac{2^{2k+2}-1}{(2k+2)!}B_{k+1}\,\theta^{2k+1},$$
where $B_k$ is Bernoulli's number, we obtain
$$\no\label{b1} 1^{2k}-2^{2k}+3^{2k}-\dots = 0 \quad (k=1,2,\dots), \\
  \no\label{b2} 1^{2k+1}-2^{2k+1}+\dots = (-1)^k\frac{2^{2k+2}-1}{2k+2}B_{k+1} \quad (k=0,1,\dots), \\
  \no\label{b3} 1^{2k+1}-3^{2k+1}+5^{2k+1}-\dots = 0 \quad (k=0,1,\dots).$$
Similarly, starting from
$$\no e^{i\theta}-e^{3i\theta}+e^{5i\theta}-\dots = \frac{e^{i\theta}}{1+e^{2i\theta}} = \frac12\sec\theta,$$
and remembering that
$$\sec\theta = 1+\Sum_1^\infty \frac{E_k\,\theta^{2k}}{(2k)!},$$
where $E_k$ is Euler's number, we obtain \eqref{b3} and also
$$\no 1^{2k}-3^{2k}+5^{2k}-\dots = \frac12(-1)^kE_k \quad (k=1,2,\dots).$$
We observe in passing that \eqref{b2}, for $k=0$, is
$$\no\label{quarter} 1-2+3-4+\dots = \frac14,$$
which is also the result of squaring \eqref{grandi} by Cauchy's rule
$$(1-1+1-\dots)(1-1+1-\dots) = 1\,.\,1-(1\,.\,1+1\,.\,1)+(1\,.\,1+1\,.\,1+1\,.\,1)-\dots.$$

(3) If we integrate \eqref{cos-alt} from $\theta=0$ to $\theta=\phi$, and
then write $\theta$ again for $\phi$, we obtain
$$\no\label{first-int} \sin\theta-\frac12\sin2\theta+\frac13\sin3\theta-\dots = \frac12\theta \quad (-\pi<\theta<\pi).$$
This series is convergent. A second integration gives
$$\no\label{second-int} 1-\cos\theta-\frac{1-\cos2\theta}{2^2}+\frac{1-\cos3\theta}{3^2}-\dots = \frac14\theta^2.$$
Here we may include the limits,\footnote{The series being uniformly
convergent for all $\theta$.} and $\theta=\pi$ gives
$$\no 1+\tfrac{1}{3^2}+\tfrac{1}{5^2}+\dots = \frac18\pi^2.$$
Since
$$1+\tfrac1{3^2}+\tfrac1{5^2}+\dots = 1+\tfrac1{2^2}+\tfrac1{3^2}+\dots-\tfrac1{2^2}-\tfrac1{4^2}-\dots = \left(1-\tfrac14\right)\left(1+\tfrac1{2^2}+\tfrac1{3^2}+\dots\right),$$
we deduce
$$\no 1+\tfrac1{2^2}+\tfrac1{3^2}+\dots = \frac16\pi^2, \no 1-\tfrac1{2^2}+\tfrac1{3^2}-\dots = \frac1{12}\pi^2,$$
and so
$$\no\label{last-int} \cos\theta-\frac{\cos2\theta}{2^2}+\frac{\cos3\theta}{3^2}-\dots = \frac1{12}\pi^2-\frac14\theta^2 \quad (-\pi\le\theta\le\pi).$$

Further integrations lead to the summation of
$\sum(-1)^{n-1}n^{-2k}\cos n\theta$ and $\sum(-1)^{n-1}n^{-2k-1}\sin n\theta$
by means of the Bernoullian functions.

(4) Alternatively, we could, by a more daring calculation, deduce
\eqref{second-int} from \eqref{grandi} and \eqref{b1}, arguing that
$$\Sum_1^\infty (-1)^{n-1}\frac{1-\cos n\theta}{n^2} &= \Sum_{n=1}^\infty \frac{(-1)^{n-1}}{n^2}\Sum_{k=0}^\infty (-1)^k\frac{(n\theta)^{2k+2}}{(2k+2)!} \\
  &= \Sum_{k=0}^\infty \frac{(-1)^k\theta^{2k+2}}{(2k+2)!}\Sum_{n=1}^\infty (-1)^{n-1}n^{2k} = \frac12\theta^2(1-1+1-\dots) = \frac14\theta^2.$$
Indeed we could generalize this argument. Suppose that
$$f(\theta) = a_0+a_1\theta^2+a_2\theta^4+\dots$$
is convergent for all $\theta$. Then the argument suggests that
$$\no \Sum_1^\infty (-1)^{n-1}\frac{f(n\theta)}{n^2} &= \Sum_{n=1}^\infty \frac{(-1)^{n-1}}{n^2}\Sum_{l=0}^\infty a_l(n\theta)^{2l} = \sum_{l=0}^\infty a_l\theta^{2l}\sum_{n=1}^\infty (-1)^{n-1}n^{2l-2} \\
  &= a_0\left(\frac1{1^2}-\frac1{2^2}+\frac1{3^2}-\dots\right)+a_1\theta^2(1-1+1-\dots) = \frac1{12}a_0\pi^2+\frac12a_1\theta^2.$$
This is plainly not true generally; for example, it is false when
$f(\theta)=e^{-\theta^2}$: but it is true for quite extensive classes of
functions. Thus, if $f(\theta)$ is the Bessel function
$$J_0(\theta) = 1-\frac{\theta^2}{2^2}+\frac{\theta^4}{2^2\,.\,4^2}-\dots,$$
it gives
$$\no\label{bessel} J_0(\theta)-\frac{J_0(2\theta)}{2^2}+\frac{J_0(3\theta)}{3^2}-\dots = \frac1{12}\pi^2-\frac18\theta^2 \quad (-\pi<\theta<\pi).$$

(5) From \eqref{sin} we deduce
$$\sum_1^\infty \cos n\theta\sin n\phi = \frac14\{\cot\frac12(\phi+\theta)+\cot\frac12(\phi-\theta)\} = \dfrac12\,\frac{\sin\phi}{\cos\theta-\cos\phi},$$
and so
$$\frac{\cos m\theta-\cos m\phi}{\cos\theta-\cos\phi} = 2\Sum_{n=1}^\infty \frac{\sin n\phi}{\sin\phi}\cos n\theta(\cos m\theta-\cos m\phi)$$
for any positive integral $m$. If we integrate this equation from
$\theta=0$ to $\theta=\pi$ (ignoring any difficulties about the range of
$\theta$ over which it may be expected to be valid\footnote{We have to
expect trouble with \eqref{sin} for $\theta=0$ or $\theta=2\pi$, since the
left-hand side vanishes identically and the right-hand side has an infinity,
and here for values of $\theta$ for which $\cos\theta=\cos\phi$. But it is
not unreasonable to suppose that these difficulties will disappear when we
multiply by the factor $\cos m\theta-\cos m\phi$; and the result seems to
justify our expectation.}), we obtain
$$\no\label{int-a} \int_0^\pi \frac{\cos m\theta-\cos m\phi}{\cos\theta-\cos\phi}\,d\theta = \pi\frac{\sin m\phi}{\sin\phi},$$
which may be verified in various ways.

(6) It follows from \eqref{sin} and \eqref{sin-alt} that
$$\sin\theta+\sin3\theta+\dots = \frac12\cosec\theta, \qquad \sin2\theta+\sin4\theta+\dots = \frac12\cot\theta.$$
If we multiply these equations by $\theta$, integrate from $\theta=0$ to
$\theta=\frac12\pi$, and observe that
$$\int_0^{\frac12\pi} \theta\sin(2n+1)\theta\,d\theta = \frac{(-1)^n}{(2n+1)^2}, \qquad \int_0^{\frac12\pi} \theta\sin2n\theta\,d\theta = (-1)^{n-1}\frac{\pi}{4n},$$
we obtain
$$\no\label{int-b} \int_0^{\frac12\pi} \frac{\theta}{\sin\theta}\,d\theta = 2\left(1-\tfrac1{3^2}+\tfrac1{5^2}-\dots\right), \\
  \no\label{int-c} \int_0^{\frac12\pi} \theta\cot\theta\,d\theta = \frac12\pi\log2.$$
These formulae also may be verified independently.

\section{First definitions.} The results of the formal calculations of
\S1.2 are correct wherever they can be checked: thus all of the formulae
\eqref{first-int}--\eqref{last-int}, \eqref{bessel}, and
\eqref{int-a}--\eqref{int-c} are correct. It is natural to suppose that the
other formulae will prove to be correct, and our transformations
justifiable, if they are interpreted appropriately. We should then be able
to regard the transformations as shorthand representations of more complex
processes justifiable by the ordinary canons of analysis. It is plain that
the first step towards such an interpretation must be some definition, or
definitions, of the `sum' of an infinite series, more widely applicable than
the classical definition of Cauchy.

This remark is trivial now: it does not occur to a modern mathematician
that a collection of mathematical symbols should have a `meaning' until one
has been assigned to it by definition. It was not a triviality even to the
greatest mathematicians of the eighteenth century. They had not the habit of
definition: it was not natural to them to say, in so many words, `by $X$ we
\emph{mean} $Y$'. There are reservations to be made, to which we shall
return in \SS1.6--7; but it is broadly true to say that mathematicians
before Cauchy asked not `How shall we \emph{define} $1-1+1-\dots$?' but
`What \emph{is} $1-1+1-\dots$?', and that this habit of mind led them into
unnecessary perplexities and controversies which were often really verbal.

It is easy now to pick out one cause which aggravated this tendency, and
made it harder for the older analysts to take the modern, more
`conventional', view. It generally seems that there is only one sum which it
is `reasonable' to assign to a divergent series: thus all `natural'
calculations with the series \eqref{1-1} seem to point to the conclusion
that its sum should be taken to be $\frac12$. We can devise arguments
leading to a different value,\footnote{See \S1.6(2).} but it always seems
as if, when we use them, we are somehow `not playing the game'.

The reason for this is fairly obvious. The simplest argument for
\eqref{grandi} is `$s=1-1+1-\dots=1-(1-1+1-\dots)=1-s$, and so
$s=\frac12$': we thus obtain the value $\frac12$, whatever our definition,
provided only that it satisfies certain very natural conditions.

Let us suppose, for example, that we have given any definition of the sum
of a series which satisfies the following axioms:

\begin{items}
\item[($A$)] \emph{if}\enskip$\sum a_n=s$\enskip\emph{then}\enskip$\sum ka_n=ks$;
\item[($B$)] \emph{if}\enskip$\sum a_n=s$\enskip\emph{and}\enskip$\sum b_n=t$,\enskip
\emph{then}\enskip$\sum(a_n+b_n)=s+t$;
\item[($C$)] \emph{if}\enskip$a_0+a_1+a_2+\dots=s$\enskip\emph{then}\enskip
$a_1+a_2+a_3+\dots=s-a_0$, \emph{and conversely}.
\end{items}

\noindent Actually, all definitions which we shall use satisfy ($A$) and
($B$), and most, though not all, satisfy ($C$). Then, if $1-1+1-\dots=s$,
we have
$$s = 1-1+\dots = 1+(-1+1-\dots) = 1-(1-1+\dots) = 1-s:$$
here we have used only ($A$) and ($C$). Similarly, if $1-2+3-4+\dots=s$, we
have
$$s = 1-2+3-\dots &= 1+(-2+3-4+\dots) = 1-(2-3+4-\dots) \\
  &= 1-(1-1+1-\dots)-(1-2+3-\dots) = 1-\frac12-s,$$
and so $s=\frac14$, in agreement with \eqref{quarter}. Here we have used all
of ($A$), ($B$), and ($C$).
